\begin{figure}[!bp]
\centering
\begin{tikzpicture}[x=0.88cm,y=0.75cm,font=\small]
\draw[-{Stealth}] (-1.3,0)--(11.5,0) node[right] {$\mu(p)$};
\draw[densely dotted,gray] (3.8,0.2)--(3.8,3.45);
\draw[densely dotted,gray] (6.2,0.2)--(6.2,3.45);
\draw[blue!65!black,line width=1pt] (-1,3)--(5,3);
\draw[blue!65!black,line width=4pt] (0,3)--(3.8,3);
\draw[teal!75!black,line width=1pt] (2.4,2)--(7.6,2);
\draw[teal!75!black,line width=4pt] (3.8,2)--(6.2,2);
\draw[violet!75!black,line width=1pt] (5,1)--(11,1);
\draw[violet!75!black,line width=4pt] (6.2,1)--(10,1);
\draw[blue!65!black,fill=white] (-1,3) circle (2pt);
\draw[blue!65!black,fill=white] (5,3) circle (2pt);
\draw[teal!75!black,fill=white] (2.4,2) circle (2pt);
\draw[teal!75!black,fill=white] (7.6,2) circle (2pt);
\draw[violet!75!black,fill=white] (5,1) circle (2pt);
\draw[violet!75!black,fill=white] (11,1) circle (2pt);
\fill[blue!65!black] (2,3) circle (2pt);
\fill[teal!75!black] (5,2) circle (2pt);
\fill[violet!75!black] (8,1) circle (2pt);
\node[above] at (2,3.12) {$I_{t_1}$};
\node[above] at (5,2.12) {$I_{t_2}$};
\node[above] at (8,1.12) {$I_{t_3}$};
\draw (0,0.08)--(0,-0.08) node[below] {$a=q_0$};
\draw (2,0.08)--(2,-0.08) node[below] {$t_1$};
\draw (3.8,0.08)--(3.8,-0.08) node[below] {$q_1$};
\draw (5,0.08)--(5,-0.08) node[below] {$t_2$};
\draw (6.2,0.08)--(6.2,-0.08) node[below] {$q_2$};
\draw (8,0.08)--(8,-0.08) node[below] {$t_3$};
\draw (10,0.08)--(10,-0.08) node[below] {$b=q_3$};
\end{tikzpicture}
\caption{Finite gluing, schematically with three selected families.
The open intervals $I_{t_j}$ are the neighborhoods where the local
families approximate correctly. Each thick closed segment
$[q_{j-1},q_j]$ lies strictly inside its corresponding open interval.
For a label from family $j$ that stays active along a convergent sequence,
the common shifts in \eqref{eq:affine-shifts} force the limiting exponent
minimum $\mu(p)$ into that closed segment. Either adjacent family remains
valid when $\mu(p)$ equals a cut.}
\label{fig:interval-gluing}
\end{figure}
